% Section 8 — Threats to validity
% Page budget: ~1.5 pages in acmart sigconf two-column (~1200-1400 words).
% Standard ICSE structure: construct / internal / external / statistical-conclusion.
% Source-of-truth for honest caveats: DOCS/docs8/PAPER-FINDINGS.md
%                                     "CRITICAL framing reminders" section.
% This is the section where the pre-paper audit's honest disclosures land.

\section{Threats to Validity}
\label{sec:threats}

We discuss four categories of threat following the standard organisation: construct, internal, external, and statistical-conclusion validity. Each category names the threats we have identified, the steps taken to mitigate them, and the residual risks that bear on how the paper's claims should be read.

\subsection{Construct validity}
\label{sec:threats:construct}

\paragraphhead{Hit@K assumes serial candidate review.} The Hit@1, Hit@5, and reciprocal-rank metrics implicitly model an engineer who reviews retrieved candidates serially and adopts the first useful one. In practice an experienced engineer may scan all five suggestions in parallel and disambiguate by comparing them, so the serial-review interpretation is a conservative estimate of how the agent's output is consumed. Reciprocal rank is a partial mitigation in that it credits the system for ranking the right ticket high regardless of consumption pattern, but it does not capture the parallel-scan workflow directly.

\paragraphhead{The triage label is taxonomy-authored, not engineer-adjudicated.} Per-window triage labels (\emph{ticket-worthy}, \emph{borderline}, \emph{noise}) are produced by the scenario taxonomy that controls fault injection rather than by an engineer reviewing the resulting telemetry. The labels reflect the \emph{intended} severity rather than an engineer-adjudicated one. The label-leakage discipline (\S\ref{sec:threats:internal}) and the chronological visibility rule (\S\ref{sec:evaluation:splits}) keep the model from recovering the taxonomy through telemetry shortcuts, but the labels themselves remain construct-authored.

\paragraphhead{Dollar-equivalent cost depends on a fixed per-skill cost table.} The cost-savings numbers in \S\ref{sec:results:cost} are computed against a specific per-skill cost table whose assumptions are documented in Appendix~A. The language-model verifier's cost dominates the dollar axis at $\$0.0005$ per call at gpt-4o-mini token rates, sourced from prior telemetry measurement. Different language-model providers would shift the absolute dollar axis proportionally. The savings \emph{percentage} is invariant to that shift because the always-on counterfactual and the agent's actual cost scale together, but the percentage figure assumes the cost-table relative ratios---in particular that the verifier is roughly two orders of magnitude more expensive than the cheap-path primitives---are representative. A configuration in which the language-model verifier is replaced by a self-hosted small model would compress the cost gap and reduce the savings figure.

\paragraphhead{The pages-per-incident metric is evaluated on test splits without natural multi-window outages.} The page-suppression rule is exercised on the test splits as designed---$20$, $1$, and $18$ suppression events fire across \dsob{}, \dsotel{}, and \dswol{} respectively---but the splits do not oversample the multi-hour multi-window outages that the suppression rule is most useful on. The pages-per-incident figure of $1.000$ on all three datasets is what the operational target requires; it should be read as ``the rule does not produce duplicate pages on test windows that do not naturally cluster'' rather than as a stress test of the rule under heavy multi-window load.

\subsection{Internal validity}
\label{sec:threats:internal}

\paragraphhead{Label leakage.} The primary internal-validity concern in a synthetic-but-labelled dataset is the possibility that the telemetry surface inadvertently encodes the scenario taxonomy the model is supposed to learn from the joint distribution. We mitigate this with a discipline applied at dataset-construction time and enforced by code review: no telemetry field may name our scenario taxonomy, no log key may carry a fault-type string, no span attribute may identify the fault-injection run, and no alert name may reference a family the engineer would have to discover from symptoms. The discipline applies to all three datasets we construct in-lab; \dswol{} draws on existing real Apache Jira tickets and therefore inherits whatever labelling discipline the Apache projects' bug trackers follow.

\paragraphhead{Out-of-fold training for the L1 stacker.} The numeric-feature classifier and the L1 triage stacker are fit via out-of-fold cross-validation so that every training-fold window is scored by a model that did not observe it during fit. Headline in-distribution numbers are computed against out-of-fold predictions.

\paragraphhead{Determinism and reproducibility.} Language-model calls run at temperature zero with a fixed seed; bootstrap confidence intervals use seed $42$ across all pipelines so that paired differences are computed on identical resamples (\S\ref{sec:evaluation:stats}). Cache, trace, and per-skill prediction outputs are persisted under a deterministic directory layout so that any specific window's decision can be reproduced by re-running a single command (\S\ref{sec:evaluation:impl}).

\paragraphhead{The structural verifier skip on \dswol{} relies on cited prior measurement.} The agent refuses to invoke the language-model verifier on \dswol{} because the dataset's profile lists it as known-harmful, citing prior measurement that the verifier degrades Hit@5 on real Apache Jira tickets. We do \emph{not} re-validate that prior measurement in this paper. A reviewer concerned with the structural skip's empirical basis should consult the prior work that produced the original measurement; the agent's design exposes the verifier-calibration table as a configurable entry, so the structural skip can be re-evaluated against new datasets without code changes.

\subsection{External validity}
\label{sec:threats:external}

\paragraphhead{Three datasets are not a population sample.} The three datasets used in this paper---a controlled synthetic e-commerce telemetry corpus, a structurally different polyglot synthetic application, and a real Apache Jira ticket corpus---are not independent samples from a population of incident-triage workloads. \dsob{} and \dsotel{} share the same dataset-construction discipline; \dswol{} replaces the telemetry side with engineer-voice prose rather than offering a third synthetic application. Generalisation claims for the agent's design should be read as well-supported within the regime of telemetry-paired-with-Jira-style-text, and as future-work claims for qualitatively different operational regimes (cloud-native serverless workloads, edge-deployed mobile telemetry, mainframe operations logs, security incident-response operations).

\paragraphhead{Single-cluster, single-region.} The synthetic-telemetry collection runs on a single Kubernetes cluster in a single region. The agent has therefore never been exposed to multi-region failure modes (cross-region latency, partial-region outage, control-plane partition) that real production fleets contend with. The cross-dataset transfer to \dsotel{} adds Kafka-based asynchronous communication as a structurally novel dependency boundary, but the multi-region question remains untested.

\paragraphhead{Synthetic Jira corpus on the two telemetry datasets.} \dsob{} and \dsotel{} pair telemetry with a Jira corpus generated by a humaniser rather than written by engineers, so it carries the humaniser's lexical fingerprint. We anchor the humaniser's length, code-block, and comment-density distributions to the real-Jira distribution from the public TAWOS corpus to limit the divergence. The corresponding pairing on \dswol{} is the inverse: real engineer-voice Jira prose paired with telemetry-style query text we extract from the ticket's symptom paragraph. Neither pairing perfectly models a production deployment in which an organisation owns both its telemetry and its ticket text.

\paragraphhead{The \dswol{} evaluable subset is small.} The family-held-out test split on \dswol{} produces $304$ test windows but only $55$ with a time-visible gold ticket under the visibility rule (\S\ref{sec:evaluation:datasets:wol}). Confidence intervals on \dswol{} numbers are correspondingly wider than on the synthetic datasets. The headline ordering (\dswol{} produces the highest point estimate on all four headline metrics) is stable at the lower CI bound, but the absolute magnitude of the cross-dataset comparison is more uncertain than the point estimates alone suggest.

\paragraphhead{No L1-retrained variant on \dsotel{}.} The cross-application transfer to \dsotel{} is reported as zero-shot only; we do not refit the L1 numeric-feature classifier on \dsotel{}'s training split. A team deploying the agent on a new application would typically refit the small L1 component (a fast linear stacker) on per-application data; the L1-retrained variant would tighten the cross-application comparison but is left to follow-up.

\paragraphhead{The plan-diversity number is contingent on the dataset's window-type taxonomy.} The controller emits four of six possible plan identifiers on the telemetry datasets and one of six on \dswol{}, because real Jira tickets do not carry the fault-window taxonomy that drives four of the six controller branches. The plan-diversity metric is therefore better read as ``the controller branches when the dataset gives it a reason to branch'' than as a universal property of the controller.

\subsection{Statistical-conclusion validity}
\label{sec:threats:statistical}

\paragraphhead{The paired-bootstrap procedure assumes resamples are exchangeable.} Confidence intervals reported in \S\ref{sec:results}--\S\ref{sec:cross-app} are produced by $1{,}000$-resample paired bootstrap at seed $42$ (\S\ref{sec:evaluation:stats}). The procedure assumes that test-set windows are exchangeable for the purpose of resampling, which is true under the i.i.d.-within-split construction we use, but would be violated by datasets with strong within-split temporal autocorrelation. Our datasets are family-stratified and chronologically separated across splits, which limits within-split autocorrelation but does not eliminate it entirely.

\paragraphhead{Most paired-delta intervals collapse to a point.} Of approximately eighty ablation cells we ran across the three datasets (\S\ref{sec:results:ablation},~\S\ref{sec:results:capmask}), the paired-delta interval is exactly $[0, 0]$ on most cells because every test-set window scores identically across the ablation arm and the baseline arm. This is informative---it tells us the ablated primitive does not change rank-one ranking on a single window in the test split---but it also means most non-significant findings have not been meaningfully tested for power. The right reading of a $[0, 0]$ interval at $p = 1.0$ is ``the ablation has no measurable effect on this split'' rather than ``the ablation has no effect''; a larger test split could in principle surface a small effect that exactly cancels at this sample size.

\paragraphhead{The five-percent significance level for the verifier-side hypotheses.} The two ablation cells that produce significance at $p = 0.002$ (extended-trace tool harm on \dsob{}) and $p < 10^{-4}$ (numeric-feature classifier essentiality on \dsob{} and \dsotel{}) are highly significant under any reasonable correction for multiple comparisons over the eighty cells. The non-significant cells reported as such (\texttt{no\_react} at $p = 0.262$, peer-incident tool at $p = 0.180$) sit well above the conventional $0.05$ threshold and are reported as non-significant without further correction. We have not applied a Bonferroni or Benjamini--Hochberg correction because the negative claims are reported as ``not statistically distinguishable from zero,'' which is robust to multiple-comparisons concerns; a positive claim that depended on $p < 0.05$ would need the correction, but we do not make such a claim from these cells.

\paragraphhead{Effect-size reporting alongside $p$-values.} Throughout \S\ref{sec:results}--\S\ref{sec:cross-app} we report point estimates with confidence intervals alongside $p$-values rather than $p$-values alone. The headline cost-savings figure ($64.1\%$) is reported without a $p$-value because it is computed against a fixed cost-table counterfactual rather than against a stochastic baseline; an effect-size-only reading is appropriate for that metric and we make no claim about its statistical significance.
